\section{结论}

本文建立一维轴对称径向热湿耦合模型，采用半控制体 FVM 与隐式时间推进求解，
回答了四个问题，并通过解析基准、收敛性、守恒、退化、灵敏度与稳健性六类检验验证模型。

\subsection{四问答案}

\begin{table}[H]
  \centering
  \caption{四个问题的答案}
  \label{tab:answers}
  \footnotesize
  \begin{tabularx}{\textwidth}{@{}clX@{}}
    \toprule
    问题 & 答案 & 关键说明 \\
    \midrule
    1 & 1800 s：中心 $33.58\ ^\circ$C、表面 $36.79\ ^\circ$C
      & 干燥只发生在表层（外侧约 0.7 cm）；中心含水率 30 min 内完全未变 \\
    2 & 3 h：中心 $49.85\ ^\circ$C、表面 $49.97\ ^\circ$C；
        中心 $1.77$、表面 $1.01$ kg/kg
      & 双向耦合使 1800 s 中心温度比问题1 低 $1.386\ ^\circ$C \\
    3 & $\tstar=\mathbf{57.49}$ h（$=206959.9521$ s $=2.40$ 天）
      & 判据必须是全域最大值；用面积平均会低估 $37.3\%$ \\
    4 & $\tdry=\mathbf{51.0906}$ h，终态半径 $1.2000$ cm
      & $\Delta t=1$ s 收敛口径；与另一套独立实现相差 $0.2\%$ \\
    \bottomrule
  \end{tabularx}
\end{table}

\subsection{检验结论与误差预算}

Bessel 解析基准误差 $<0.001\ ^\circ$C（相对 $<0.002\%$）；空间二阶
（实测 $+2.20$）、时间不低于一阶（\textbf{不声称二阶}）；
内部界面通量抵消达机器精度 $1.782\times10^{-16}$；五项退化与一致性检验通过；
跨实现逐行移植复现差 $+11$ s（$0.004\%$）。
几何—密度一致性诊断提示密度定义、数据误差、几何假设或干物质守恒的
联合使用需核查（\textbf{条件性结论，不下``数据有误''的判断}）。

误差预算（\S\ref{sec:errbudget}）的落点是\textbf{口径主导，不是数据主导}：
\textbf{平台窗口取法一项即达 $1090$ s，比最大的数据扰动（$135$ s）还大 8 倍}，
也远大于全部数值离散误差（约 $19$ s）。

\subsection{灵敏度与稳健性}

敏感度排序 $D\gg h_m\gg h$ 在五档扰动下均不变；
但 $|S_D|$ 在 $\pm5\%$ 处为 $0.886$、在 $-50\%$ 处为 $1.794$——
\textbf{灵敏度系数是局部量，报告时必须注明档位}。
数据扰动引起的散布为 $0.22\sim67.6$ s，\textbf{全部经由``长时平台均值''
单一渠道传递}（$r=-0.9997$）。
三种区间（情景 / 参数扰动分位 / 数据扰动分位）严格分开命名，
\textbf{统计置信区间不给}。

\subsection{模型局限与推广}

\textbf{局限}：一维轴对称简化，端面效应以二维对照验证而非全面 CFD；
M0 不含蒸发潜热（M1 对照显示 3 h 中心温度低 $18.06$ K、$\tstar$ 延长
$2.84$ h，属\textbf{条件性扩展，不改答案}）；附件1 只覆盖 0--4 h 需外推；
缺少内部位移、水活度、目标成分动力学等数据。

\textbf{推广}：烘房—物料双向耦合系统模型（需空气干质量、装料量、
新风/排风质量流量）；品质后处理（需水活度关系、成分动力学、菌种与初始菌量）；
结构力学（需应变、本构、强度参数）。
